\addchap{Preface}

This book began with a list of textbooks that should exist but do not. One entry observed that many exact predictions of high-dimensional statistics and machine-learning theory come from statistical physics, from replicas, cavities, TAP equations, Kac--Rice counts and dynamical mean-field theory, and that the two communities who own them cannot read each other's arguments. Mathematicians do not read the physical derivations, and physicists do not see the intuition inside the rigorous proofs. Each side has excellent books, but there was no parallel text. The entry proposed one: the physical derivation on one side and the corresponding rigorous route on the other, with the replica method set against Guerra and Talagrand's interpolation, the cavity method against approximate message passing and its state evolution, Gaussian comparison against the convex Gaussian min--max theorem, and dynamical mean-field theory against the rigorous high-dimensional limits of stochastic gradient descent; a board in every chapter recording which predictions are proved, which are conjectures, and which are known to be wrong; and examples drawn only from inference and learning.

The resulting book has fifteen chapters in four parts. Part I introduces the two languages on the simplest disordered systems: free entropies and self-averaging, the random energy model, and the Sherrington--Kirkpatrick model. Part II is Bayesian inference, where the posterior is the Gibbs measure that physics studies and replica symmetry becomes a theorem: the Nishimori line, spiked matrices, message passing, and generalized linear models with their hard phase. Part III is convex learning, where Gordon's comparison inequality replaces the replicas: the capacity of the perceptron, ridge regression, the LASSO, logistic regression and max-margin classifiers, and kernels with the Gaussian equivalence principle. Part IV asks where algorithms go: the critical points of random landscapes, dynamical mean-field theory, and online stochastic gradient descent. Every chapter states its prediction, derives it as a physicist would, sets the derivation and the proof side by side step by step in a section called the stone, proves the central result completely in general or in a special case that keeps the idea, measures it in a short script, and ends with a status board and an account of where the translation fails.

The intended readers are graduate students and researchers in probability, statistics, machine learning and statistical physics who have met an exact prediction in a paper and want to know where it came from, why it is believed, and what has been proved. The prerequisites are linear algebra, multivariable calculus, probability at the level of a first graduate course, and some convex optimization. No spin-glass theory is assumed, and the random-matrix facts that the book uses are stated in the appendix, which also proves the Gaussian tools.

Two cautions belong at the front. First, the status of these results changes quickly. Several long-standing predictions were settled in preprints posted in 2025 and 2026, some of them, as their authors state, with proofs written or checked with the help of AI systems. The book reports such results as their authors state them, marks them as preprints, and does not count them as proved. Second, a parallel text inevitably simplifies both sides: the physical derivations are given in their cleanest form rather than in their historical one, and many proofs are given in a special case that keeps the idea, with the general theorem stated and its source named.

\paragraph{A note on authorship.} This book was written by Claude, an AI model made by Anthropic, and conceived and edited by Xiaoyu Li, who proposed the subject, its form as a parallel text, its status boards, and its examples. Every theorem is stated with its hypotheses and is proved, sketched, or attributed to the source that proves it; every physical derivation is marked as one; several results and proofs are the book's own and are marked as such. Every claim attributed to a paper was checked against the paper or a reliable account of it wherever these could be reached, every number from the book's scripts comes from code that the reader can run, and the colophon describes how the book was made. Readers who find an error are invited to report it, with the computation that fails.
